\section*{Chapter \ref{ch:b3:stable-isotopes} --- Stable-Isotope Geochemistry}

\begin{solution}{exo:b3:stable-isotopes:1}
$\alpha = 1030.5/997.9 = 1.032669$; $1000\ln\alpha = 32.15$; $\Delta = 30.5 -
(-2.1) = \qty{32.6}{\permil}$. They differ by \qty{0.45}{\permil}, the
$(\delta_A^2 - \delta_B^2)/2000 = 0.46$ of the proposition.
\end{solution}

\begin{solution}{exo:b3:stable-isotopes:2}
$\delta^{18}\mathrm O = 0.69\times(-25) - 13.6 = \qty{-30.9}{\permil}$;
$\delta\mathrm D = 5.6\times(-25) - 100 = \qty{-240}{\permil}$; $d = -240 +
8\times30.85 = \qty{6.8}{\permil}$.
\end{solution}

\begin{solution}{exo:b3:stable-isotopes:3}
$\Delta_B = -26.0 - 1.0 = \qty{-27.0}{\permil}$; $f_{\mathrm{org}} = (-5 -
1.0)/(-27.0) = 0.22$.
\end{solution}

\begin{solution}{exo:b3:stable-isotopes:4}
$T^2 = 0.0391\times10^6/(0.600 - 0.154) = \num{87668}$, $T = \qty{296.1}{K}$,
\qty{22.9}{\celsius}.
\end{solution}

\begin{solution}{exo:b3:stable-isotopes:5}
$f^{0.01} = 967/987$, $f = \exp(100\ln0.97974) = 0.129$: \qty{87}{\percent} of
the vapour has rained out. The rain: $1.0100\times967 - 1000 =
\qty{-23.3}{\permil}$.
\end{solution}

\begin{solution}{exo:b3:stable-isotopes:6}
$\alpha = 1000/960 = 1.041667$; $1000\ln\alpha = 40.82$; $1000(\alpha - 1) =
41.67$; $\Delta = 40$. Both approximations miss by about \qty{0.8}{\permil}, the
second by twice as much as the first: at \qty{40}{\permil} the expansion to
first order is no longer enough.
\end{solution}

\begin{solution}{exo:b3:stable-isotopes:7}
$\delta_c = (1.50 + 29.99)/0.97001 = \qty{32.46}{\permil}$ VSMOW; $1000\ln\alpha
= 1000\ln1.032464 = 31.95$; $T = 18030/(31.95 + 32.42) = \qty{280.1}{K}$,
\qty{7.0}{\celsius}. The fractionation changes by about \qty{0.23}{\permil} per
kelvin, so \qty{0.1}{\permil} on the water is about \qty{0.4}{\celsius}.
\end{solution}

\begin{solution}{exo:b3:stable-isotopes:8}
Ice volume $125\times0.011 = \qty{1.375}{\permil}$; the rest, $1.77 - 1.375 =
\qty{0.40}{\permil}$. Near \qty{2}{\celsius} the slope is $18030/275.15^2 =
\qty{0.238}{\permil.K^{-1}}$, so the deep water was about \qty{1.7}{\celsius}
colder.
\end{solution}

\begin{solution}{exo:b3:stable-isotopes:9}
$\Delta^{33}\mathrm S = 8.2 - 0.515\times10.0 = \qty{3.05}{\permil}$, far from 0:
mass-independent sulfur, which forms only without an ozone screen. The pyrite
formed under an atmosphere without significant oxygen, before about
\qty{2.45}{\Ga}.
\end{solution}

\begin{solution}{exo:b3:stable-isotopes:10}
The heavy isotope is conserved in the closed system: $R_0 = fR_v + (1 -
f)\alpha R_v$, so $R_v/R_0 = 1/[f + \alpha(1 - f)]$. For $f = 0.2$, $\alpha =
1.01$: $1/1.008$, the vapour is \qty{7.9}{\permil} lighter; Rayleigh gives
$0.2^{0.01}$, \qty{16.0}{\permil} lighter. Removing each condensate at once
drives the vapour twice as far.
\end{solution}

\begin{solution}{exo:b3:stable-isotopes:11}
$M = m(\delta_s - \Delta)/\Delta = 2000\times(-60 + 2.5)/(-2.5) = \qty{46000}{Gt}$.
Then $m = 46000\times(-2.5)/(-25 + 2.5) = \qty{5100}{Gt}$, a little below the
published 6000 to \qty{12000}{Gt}, whose upper end corresponds to the larger
excursion.
\end{solution}

\begin{solution}{exo:b3:stable-isotopes:12}
$\delta_v = (1000 + \delta_{v,0})\exp\int(\alpha - 1)\,\mathrm d\ln f - 1000$, so
$\mathrm d\delta_v = (1000 + \delta_v)(\alpha - 1)\,\mathrm d\ln f$, and $f =
e_s/e_{s,0}$ gives $\mathrm d\ln f = \mathrm d\ln e_s$. At \qty{-20}{\celsius},
$\alpha = 1.01872$: $0.01872\times950.9\times0.0960 =
\qty{1.71}{\permil.K^{-1}}$ for the vapour.
\end{solution}

\begin{solution}{pb:b3:stable-isotopes:1}
\textbf{1.} $\alpha^{18}$: 1.01605, 1.01691, 1.01779, 1.01872, 1.01968,
1.02068; $\alpha^{\mathrm D}$: 1.1424, 1.1524, 1.1630, 1.1744, 1.1866, 1.1997
(at $-5$ to \qty{-30}{\celsius}).

\textbf{2.} $980\times0.425^{0.01605} - 1000 = \qty{-33.4}{\permil}$.

\textbf{3.} $1.01691\times966.6 - 1000 = \qty{-17.0}{\permil}$.

\textbf{4.} Vapour $966.6\times(0.3976)^{0.01779} - 1000 = \qty{-49.1}{\permil}$;
snow $1.01872\times950.9 - 1000 = \qty{-31.3}{\permil}$.

\textbf{5.} Vapour $950.9\times(0.3680)^{0.01968} - 1000 = \qty{-67.6}{\permil}$;
snow $1.02068\times932.4 - 1000 = \qty{-48.3}{\permil}$.

\textbf{6.} $(932.4/980)\times0.0622 = 0.059$: \qty{5.9}{\percent} of the
oxygen-18, against \qty{6.2}{\percent} of the water: the vapour has lost the
heavy isotope faster.

\textbf{7.} $(-17.0 + 48.3)/20 = \qty{1.57}{\permil}$ per degree.

\textbf{8.} $t = (-48.3 + 13.6)/0.69 = \qty{-50}{\celsius}$. The relation is
for surface temperatures; over an ice sheet the surface is much colder than the
air where the snow forms, and the model ignores mixing and kinetic effects.

\textbf{9.} Vapour $-247.5$, $-352.5$, \qty{-462.7}{\permil}; snow $\delta\mathrm D
= -132.8$, $-239.6$, \qty{-355.4}{\permil}.

\textbf{10.} $d = -132.8 + 8\times17.02 = 3.3$; then 10.7 and
\qty{31.3}{\permil}: the excess grows inland. The local slope of the
distillation, $\mathrm d\delta\mathrm D/\mathrm d\delta^{18}\mathrm O = (\alpha^{\mathrm D} -
1)(1000 + \delta\mathrm D)/[(\alpha^{18} - 1)(1000 + \delta^{18}\mathrm O)]$, falls below
8 as the vapour grows very light ($1000 + \delta\mathrm D$ shrinks much faster than
$1000 + \delta^{18}\mathrm O$), so the snow drifts above the line.

\textbf{11.} The line gives $8\times(-48.34) + 10 = \qty{-376.7}{\permil}$; the
snow is \qty{21}{\permil} above it. The model has only equilibrium
fractionation; real snow forms in supersaturated air, with a kinetic effect,
and air masses mix with others on the way.

\textbf{12.} $\delta_c = (3.25 + 29.99)/0.97001 = 34.27$; $1000\ln\alpha =
33.69$; $T = 18030/66.11 = \qty{272.7}{K}$, \qty{-0.4}{\celsius}.

\textbf{13.} $\delta_c = 36.09$, $1000\ln\alpha = 35.46$, $T = \qty{265.6}{K}$,
\qty{-7.5}{\celsius}: colder than seawater can be without freezing. The
change cannot be all temperature.

\textbf{14.} $\delta_w = 125\times0.011 = \qty{+1.375}{\permil}$; $1000\ln(1036.09/1001.375)
= 34.08$; $T = 18030/66.50 = \qty{271.1}{K}$, \qty{-2.0}{\celsius}: a cooling of
\qty{1.6}{\celsius}.

\textbf{15.} $\Delta_{47} = 0.0391\times10^6/272.71^2 + 0.154 = 0.6797$ and
$0.0391\times10^6/271.12^2 + 0.154 = 0.6859$: a difference of
\qty{0.006}{\permil}.

\textbf{16.} It measures how the heavy isotopes are arranged within the
mineral, an internal equilibrium that depends on temperature only.

\textbf{17.} $f_{\mathrm{org}} = (-5 - 0.5)/(-28) = 0.20$.

\textbf{18.} $m = 46000\times(-3.0)/(-60 + 3.0) = \qty{2400}{Gt}$.

\textbf{19.} $46000\times(-3.0)/(-25 + 3.0) = \qty{6300}{Gt}$.

\textbf{20.} $46000\times(-3.0)/(-5 + 3.0) = \qty{69000}{Gt}$, more than the
reservoir itself, in a geologically brief time: implausible.

\textbf{21.} The depth and extent of carbonate dissolution in the core, which
measures how much carbon entered the ocean, independently of its composition;
with the excursion's size it fixes $\delta_s$.

\textbf{22.} The snow at Dome Ilse has $\delta^{18}\mathrm O =
\textbf{\qty[text-series-to-math=true]{-48.3}{\permil}} $ (\qty{-48.4}{\permil}
with $\alpha$ varying continuously).
\end{solution}
